\documentclass[12pt]{article}
\usepackage[margin=1in]{geometry}
\usepackage{amsmath}
\usepackage{enumitem}
\usepackage{hyperref}

\begin{document}

\begin{center}
\textbf{SOSC 13200 (W26): Social Science Inquiry II --- Midterm Practice Sheet}
\end{center}

\noindent\textbf{Time:} 80 minutes \\
\textbf{Notes:} Closed notes. No R / no computers. \\
\textbf{Allowed:} Pencil/pen; non-programmable calculator. \\
\textbf{Instructions:} Answer all questions. Show your work for calculations; partial credit is available. If you use notation, define it.

\vspace{0.5em}
\noindent\textbf{Provided formulas (you may use these)}
\begin{itemize}
  \item Independence: $\Pr[A \cap B]=\Pr[A]\Pr[B]$
  \item Conditional probability: $\Pr[A\mid B]=\dfrac{\Pr[A\cap B]}{\Pr[B]}$ (if $\Pr[B]>0$)
  \item Bayes' Rule: $\Pr[A\mid B]=\dfrac{\Pr[B\mid A]\Pr[A]}{\Pr[B]}$
  \item Law of total probability (binary $A$): $\Pr[B]=\Pr[B\mid A]\Pr[A]+\Pr[B\mid A^c]\Pr[A^c]$
  \item Wald estimator (binary instrument $Z$):
  \[
  \hat\beta_{\text{Wald}}=\frac{\bar Y_{Z=1}-\bar Y_{Z=0}}{\bar X_{Z=1}-\bar X_{Z=0}}
  \]
\end{itemize}

\vspace{0.5em}
\noindent\textbf{Practice Problem 1: Inference, measurement, and research design}
\begin{enumerate}[label=\arabic*.]
  \item \cite{holland1986statistics}: Define descriptive or associational statistics and causal inference. Give one example of each.
  \item \cite{king1994designing} Overview the elements of KKV's definition of scientific research. Explain each element in 1--2 sentences. 
  \item \cite{adcockcollier2001measurement} Provide one example of a concept that is hard to measure and describe a reasonable operationalization.
  \item \cite{angrist1991does}, \cite{cardkrueger1994minimum}: Identify the unit of analysis, study population, treatment, and outcome in one of these papers. Explain why the chosen treatment is plausibly exogenous, i.e., why we can infer a causal relationship between the treatment and outcome.
\end{enumerate}

\vspace{0.5em}
\noindent\textbf{Practice Problem 2: Causal assumptions} \cite{holland1986statistics}

\noindent The City of Chicago wants to estimate the \emph{short-term causal effect} of installing
brighter LED streetlights on nighttime pedestrian activity.

\noindent Define the outcome as $Y_i = 1$ if resident $i$ walked outside (for any reason) between
9pm and 10pm on the survey day, and $0$ otherwise. Assume everyone responds to the survey
and there are no measurement issues.

\vspace{0.5em}
\noindent For each proposed study design below, identify \textbf{one Holland assumption that is most
likely to be violated} (or is most questionable), and briefly explain why (1--2 sentences).
Choose from:

\begin{itemize}
  \item \textbf{Temporal stability}
  \item \textbf{Causal transience} 
  \item \textbf{Unit homogeneity}
  \item \textbf{Constant causal effect}
  \item \textbf{Independence}
\end{itemize}

\noindent \textbf{Answer format (required):} \emph{Assumption:} \_\_\_\_; \emph{Why:} \_\_\_\_.

\vspace{0.5em}
  \begin{enumerate}[label=\arabic*.]
    \item The city offers a website where residents can \textbf{request} brighter streetlights on their block.
    Blocks are upgraded as requests come in. Several weeks after opening the website, the city surveys residents on upgraded blocks and
    non-upgraded blocks and compares $Y$ across these groups.

    \item The city runs a pilot by upgrading \textbf{all streetlights only in the Loop}. They survey all residents in the Loop and in other neighborhoods and compare $Y$ across these areas.

    \item The city runs a pilot by \textbf{randomly upgrading streetlights block by block only in the Loop}. They survey all residents in the Loop, and compare $Y$ between blocks where lights were upgraded and where they were not. 
    They plan to use that estimate as \textbf{the predicted effect for all neighborhoods citywide}
    if the program is expanded.

    \item The city turns on the brighter LED streetlights on \textbf{all residential blocks} on the same day. They survey all residents the day before and the day after and take the difference in $Y$ as the policy effect.

    \item The city turns on brighter LED streetlights on \textbf{all residential blocks} for \textbf{one week}, then returns to
    normal lighting. They survey all residents on the first Monday of each week, and compare $Y$ across the two Mondays as the policy effect.

\end{enumerate}

\noindent\textbf{Practice Problem 3A: Sample spaces, random variables, conditional relationships}
\noindent \textbf{Setup:} A 3-person jury votes to \textbf{convict (C)} or \textbf{acquit (A)}.
Juror 1 votes to convict with probability $0.5$.
Juror 2 votes to convict with probability $0.5$, independently of Juror 1.

Juror 3 is a \emph{conformist}:
\begin{itemize}
  \item If Jurors 1 and 2 agree (both C or both A), then Juror 3 votes \emph{with them} with probability $0.9$.
  \item If Jurors 1 and 2 disagree, then Juror 3 votes C with probability $0.5$.
\end{itemize}

Let $\Omega$ be the sample space written as length-3 strings (e.g., ``CAC'') in order (Juror 1, Juror 2, Juror 3).
Let $X$ be the number of convictions (the number of C's).

Define $Y$ as:
\begin{itemize}
  \item $Y=1$ if the jury \textbf{convicts by majority rule} (at least two C votes),
  \item $Y=0$ otherwise.
\end{itemize}

\begin{enumerate}[label=\arabic*\alph*)]
  \item Write out the full sample space $\Omega$.
  \item Compute the probability of each outcome in $\Omega$ and then give the pmf of $X$:
  list $\Pr[X=0]$, $\Pr[X=1]$, $\Pr[X=2]$, $\Pr[X=3]$.
  \item Compute $\Pr[Y=1]$ and $\Pr[Y=0]$.
  \item Are the events $\{X=2\}$ and $\{Y=1\}$ independent? Justify using $\Pr(A\cap B)$ vs $\Pr(A)\Pr(B)$.
\end{enumerate}

\vspace{1em}

\noindent\textbf{Practice Problem 3B: Sample spaces, random variables, conditional relationships}

\noindent \textbf{Setup:} Three students (A, B, C) may submit an assignment on time (T) or late (L).
\begin{itemize}
  \item Student A is on time with probability $0.7$.
  \item Student B is on time with probability $0.6$.
  \item A and B are independent.
  \item Student C’s behavior depends on A and B:
  \begin{itemize}
    \item If both A and B are on time, then $\Pr(C=\text{T})=0.9$.
    \item If exactly one of A or B is on time, then $\Pr(C=\text{T})=0.5$.
    \item If both A and B are late, then $\Pr(C=\text{T})=0.1$.
  \end{itemize}
\end{itemize}

Let $\Omega$ be the sample space written as length-3 strings in order (A, B, C), e.g., ``TLT''.
Let $X$ be the number of on-time submissions (the number of T's).

Define $Y$ as:
\begin{itemize}
  \item $Y=1$ if \textbf{at least two} students are on time,
  \item $Y=0$ otherwise.
\end{itemize}

\begin{enumerate}[label=\arabic*\alph*)]
  \item Write out the full sample space $\Omega$.
  \item Give the pmf of $X$: list $\Pr[X=0]$, $\Pr[X=1]$, $\Pr[X=2]$, $\Pr[X=3]$.
  \item Compute $\Pr[Y=1]$ and $\Pr[Y=0]$.
  \item Are the events $\{\text{A is on time}\}$ and $\{Y=1\}$ independent? Briefly justify.
\end{enumerate}

% ------------------------------------------------------------

\vspace{1.5em}
\noindent\textbf{Practice Problem 4A: Conditional probability, Bayes' rule, and base rates}

\noindent \textbf{Setup (classification model):} A campaign model flags a person as a likely Democrat. 

\noindent Let $A$ be the event the person truly is a Democrat, and $F$ be the event the model flags them.

\noindent You are told:
\begin{itemize}
  \item $\Pr[A]=0.40$
  \item True positive rate: $\Pr[F\mid A]=0.80$
  \item False positive rate: $\Pr[F\mid A^c]=0.30$
\end{itemize}

\begin{enumerate}[label=\arabic*\alph*)]
  \item Compute $\Pr[F]$. Show your work.
  \item Compute $\Pr[A\mid F]$. (If flagged, what is the probability the person truly has $A$?)
  \item Compute $\Pr[A\mid F^c]$. Show your work.
  \item In 2--3 sentences: explain why $\Pr[A\mid F]$ is generally not equal to $\Pr[F\mid A]$.
  Name the fallacy this confusion relates to.
\end{enumerate}

\vspace{1em}
\noindent\textbf{Practice Problem B: Conditional probability, Bayes' rule, and base rates}

\noindent \textbf{Setup:} An automated system flags assignments as ``suspicious''.
Let $P$ be the event an assignment is truly plagiarized, and $F$ be the event it is flagged.

You are told:
\begin{itemize}
  \item $\Pr[P]=0.07$
  \item $\Pr[F\mid P]=0.85$
  \item $\Pr[F\mid P^c]=0.12$
\end{itemize}

\begin{enumerate}[label=\arabic*\alph*)]
  \item Compute $\Pr[F]$.
  \item Compute $\Pr[P\mid F]$.
  \item Compute $\Pr[P\mid F^c]$.
  \item In 2--3 sentences: what is the role of the base rate $\Pr[P]$ in (b)?
\end{enumerate}

\vspace{1em}

\noindent\textbf{Practice Problem 4C: Bayes with a rare outcome}

\noindent \textbf{Setup:} Let $V$ be the event a person will vote, and $F$ be the event the model flags them.

\noindent You are told:
\begin{itemize}
  \item $\Pr[V]=0.10$
  \item $\Pr[F\mid V]=0.90$
  \item $\Pr[F\mid V^c]=0.20$
\end{itemize}

\begin{enumerate}[label=\arabic*\alph*)]
  \item Compute $\Pr[F]$. Show your work.
  \item Compute $\Pr[V\mid F]$.
  \item Compute $\Pr[V\mid F^c]$.
  \item In 2--3 sentences: what role does the base rate $\Pr[V]$ play in your answer to (b)?
\end{enumerate}

\vspace{1em}

\noindent\textbf{Practice Problem 4C: Same task, using counts (confusion matrix)}

\noindent \textbf{Setup:} In a sample of 10{,}000 people, 6{,}000 truly vote and 4{,}000 do not.
The model flags 4{,}500 of the voters, and flags 800 of the non-voters.

\begin{enumerate}[label=\arabic*\alph*)]
  \item Compute $\Pr[F]$.
  \item Compute $\Pr[V\mid F]$.
  \item Compute $\Pr[V\mid F^c]$.
  \item In 2--3 sentences: explain why it can be misleading to report only the true positive rate.
\end{enumerate}

% ------------------------------------------------------------

\vspace{1.5em}
\noindent\textbf{Practice Problem 5A: Experiments, counterfactuals, and spillovers, with reference to \cite{miguel2004worms}}

\noindent \textbf{Setup (public health):} A ministry wants to estimate the effect of distributing
insecticide-treated bed nets on malaria incidence among children. The outcome is whether a child tests
positive for malaria after 3 months. Bed nets reduce mosquito bites, and may also reduce the local mosquito
population (potentially affecting other households).

\begin{enumerate}[label=\arabic*\alph*)]
  \item In potential-outcomes language, what does it mean to say there are treatment spillovers/externalities?
  \item Give two reasons why bed-net distribution might be randomized at the \emph{village level} rather than the household level.
  \item Assume positive spillovers so outcomes depend on neighbors' treatment. Explain why the no-interference assumption fails,
  and argue whether the simple treated-vs-control difference is attenuated.
  \item Give one design or analysis strategy that could help address spillovers (one sentence).
\end{enumerate}

\vspace{1em}

\noindent\textbf{Practice Problem 5B: Peer effects in education}

\noindent \textbf{Setup (education):} A district introduces a new reading curriculum. Some classrooms receive the curriculum and some do not.
Students in treated classrooms might share materials or strategies with friends in control classrooms in the same school.

\begin{enumerate}[label=\arabic*\alph*)]
  \item Give two reasons the district might randomize at the \emph{school} level rather than the classroom or student level.
  \item If there are positive spillovers from treated to control classrooms within the same school, what happens to a simple treated-vs-control difference?
  Explain briefly.
\end{enumerate}

% ------------------------------------------------------------

\vspace{1.5em}
\noindent\textbf{Practice Problem 6: Encouragement design, Wald estimator}

\noindent \textbf{Setup:} A clinic tests an encouragement intervention.
Let $Z=1$ if a patient is randomly assigned to receive a reminder text encouraging a flu shot, and $Z=0$ otherwise.
Let $X$ be an indicator for actually getting the flu shot, and let $Y$ be an indicator for getting the flu that season.

You are given group means:
\begin{itemize}
  \item $\bar Y_{Z=1}=0.08$, \quad $\bar Y_{Z=0}=0.10$
  \item $\bar X_{Z=1}=0.60$, \quad $\bar X_{Z=0}=0.45$
\end{itemize}

\begin{enumerate}[label=\arabic*\alph*)]
  \item Compute the Wald estimate $\hat\beta_{\text{Wald}}$ for the effect of $X$ on $Y$.
  \item Identify the reduced form and the first stage.
\end{enumerate}

\newpage
\bibliographystyle{apalike}
\nocite{butler2011politicians,angrist1991does,miguel2004worms,cardkrueger1994minimum}
\bibliography{../assets/bib}

\end{document}
